\documentclass[aspectratio=169,11pt]{beamer}

% ---- Minimalist theme ----
\usetheme{default}
\usecolortheme{default}
\setbeamertemplate{navigation symbols}{}
\setbeamertemplate{frametitle continuation}{}
\setbeamertemplate{itemize items}[circle]
\setbeamertemplate{enumerate items}[default]

% ---- Font: TeX Gyre Heros ----
\usepackage[T1]{fontenc}
\usepackage{tgheros}
\renewcommand{\familydefault}{\sfdefault}
\usepackage{microtype}

% ---- Packages ----
\usepackage{amsmath}
\usepackage{booktabs}
\usepackage{array}
\usepackage{tikz}
\usepackage{xcolor}
\usepackage{hyperref}
\usetikzlibrary{arrows.meta,positioning,calc,fit}

% ---- Maastricht University colors ----
\definecolor{umdark}{HTML}{001C3D}
\definecolor{umorange}{HTML}{E84E10}
\definecolor{umlight}{HTML}{4A90C4}
\definecolor{umgray}{HTML}{6B7280}
\definecolor{umbg}{HTML}{F8F9FA}
\definecolor{umfaint}{HTML}{E5E7EB}

% ---- Apply colors to Beamer ----
\setbeamercolor{normal text}{fg=umdark}
\setbeamercolor{frametitle}{fg=umdark}
\setbeamercolor{title}{fg=umdark}
\setbeamercolor{subtitle}{fg=umgray}
\setbeamercolor{author}{fg=umgray}
\setbeamercolor{date}{fg=umgray}
\setbeamercolor{institute}{fg=umorange}
\setbeamercolor{itemize item}{fg=umorange}
\setbeamercolor{itemize subitem}{fg=umlight}
\setbeamercolor{enumerate item}{fg=umorange}
\setbeamercolor{block title}{bg=umdark,fg=white}
\setbeamercolor{block body}{bg=umdark!5,fg=umdark}
\setbeamercolor{block title alerted}{bg=umorange,fg=white}
\setbeamercolor{block body alerted}{bg=umorange!8,fg=umdark}
\setbeamercolor{block title example}{bg=umlight,fg=white}
\setbeamercolor{block body example}{bg=umlight!8,fg=umdark}

% ---- Frametitle style ----
\setbeamertemplate{frametitle}{%
  \vspace{0.35cm}%
  \noindent\hspace*{0pt}%
  \parbox{\textwidth}{%
    {\usebeamerfont{frametitle}\usebeamercolor[fg]{frametitle}\insertframetitle}%
    \vspace{2pt}\\%
    {\color{umorange}\rule{\textwidth}{1.2pt}}%
  }%
  \vspace{-2pt}%
}
\setbeamerfont{frametitle}{size=\large,series=\bfseries}
\setbeamersize{text margin left=0.7cm,text margin right=0.7cm}

% ---- Footline ----
\setbeamertemplate{footline}{%
  \hbox{%
    \begin{beamercolorbox}[wd=\paperwidth,ht=2.2ex,dp=1ex]{footline}%
      \hspace{1em}{\scriptsize\color{umgray}Proposal Revision Summary -- Carrier Origin and Steganographic Detectability}%
      \hfill%
      {\scriptsize\color{umgray}\insertframenumber/\inserttotalframenumber}%
      \hspace{1em}%
    \end{beamercolorbox}%
  }%
}

% ---- Metadata ----
\title{Proposal Revision Summary}
\subtitle{Major Changes from Submitted Midway Proposal to Final Updated Version}
\newcommand{\teamauthors}{Abdul Moiz Akbar \;|\; Daria Gjonbalaj \;|\; Nico Muller-Spath\\Jimena Narvaez del Cid \;|\; David Wicker \;|\; Nikolas Zouros}
\author[Project 2.2 Team]{Project 2.2 Team}
\institute{Department of Advanced Computing Sciences\\Maastricht University}
\date{Project 2.2 \;|\; March 2026}

\begin{document}

% ============================================================
\begin{frame}[plain]
\vfill
\begin{center}
{\color{umorange}\rule{0.42\textwidth}{2pt}}\\[12pt]
{\LARGE\bfseries\color{umdark}\inserttitle}\\[10pt]
{\normalsize\color{umgray}\insertsubtitle}\\[16pt]
{\small\color{umdark}\teamauthors}\\[8pt]
{\small\color{umorange}\insertinstitute}\\[8pt]
{\footnotesize\color{umgray}\insertdate}\\[12pt]
{\color{umorange}\rule{0.42\textwidth}{2pt}}
\end{center}
\vfill
\end{frame}

% ============================================================
\begin{frame}{Agenda}
\begin{enumerate}
\item Why the proposal had to be revised
\item Motivation and research-question changes
\item Methodology and chosen-approach changes
\item Experiment and statistics changes
\item Related-work and reference corrections
\item What stayed constant across versions
\end{enumerate}
\end{frame}

% ============================================================
\begin{frame}{Why We Revised the Proposal}
\footnotesize
\begin{columns}[T]
\begin{column}{0.48\textwidth}
\begin{alertblock}{Problems in the submitted version}
\begin{itemize}
\item Some methods and references were weakly grounded or mismatched to the actual study plan.
\item The original design mixed domains in ways that made interpretation harder.
\item Several sections sounded over-polished and AI-like.
\item A few methodological choices were not realistic for the intended scope and implementation path.
\end{itemize}
\end{alertblock}
\end{column}
\begin{column}{0.48\textwidth}
\begin{exampleblock}{Revision goals}
\begin{itemize}
\item Keep the same core research problem.
\item Remove hallucination-prone or unsupported claims.
\item Make the design easier to justify and implement.
\item Tighten the writing into a clearer bachelor-level proposal.
\end{itemize}
\end{exampleblock}
\end{column}
\end{columns}

\vspace{0.08cm}
{\footnotesize\color{umgray}\textbf{Bottom line:} the final version is the same project, but reframed into a more defensible and implementable form.}
\end{frame}

% ============================================================
\begin{frame}{Motivation and Research Questions}
\footnotesize
\begin{columns}[T]
\begin{column}{0.49\textwidth}
\begin{block}{Original version}
\begin{itemize}
\item Stronger rhetorical emphasis on security, science, and practice.
\item More expansive wording about mixed real/synthetic traffic.
\item RQs were tied to a less stable methodology.
\end{itemize}
\end{block}
\end{column}
\begin{column}{0.49\textwidth}
\begin{exampleblock}{Updated version}
\begin{itemize}
\item Motivation now starts from steganography and steganalysis as one problem.
\item The core question is unchanged.
\item RQs are labeled by role: primary, exploratory, and verification.
\end{itemize}
\end{exampleblock}
\end{column}
\end{columns}

\vspace{0.05cm}
{\footnotesize\color{umgray}Main change: tighter framing without changing the actual problem.}
\end{frame}

% ============================================================
\begin{frame}{Chosen Approaches: Major Method Changes}
\tiny
\renewcommand{\arraystretch}{1.18}
\begin{tabular}{@{}p{3.4cm}p{4.9cm}p{4.9cm}@{}}
\toprule
\textbf{Area} & \textbf{Submitted version} & \textbf{Updated version} \\
\midrule
Carrier format & RGB PNG everywhere & Grayscale throughout; PNG for spatial, JPEG Q=95 for DCT \\
Quality control & BRISQUE gate + artifact rejection & Saturation + semantic screening; BRISQUE discussed, not used \\
Spatial embedding & PRNG-keyed pixel selection, $k=1/2$ framing & Sequential row-major LSB, with payload levels defined by fill rate \\
Frequency embedding & DCT-QIM on PNG carriers & DCT-LSB on actual JPEG carriers via quantized coefficients \\
Domain comparison & Method contrast entangled with carrier handling & Full pipeline contrast: LSB+PNG vs DCT-LSB+JPEG \\
Detector scope & SRM+EC primary baseline & Training-free primary detectors; SRNet optional \\
\bottomrule
\end{tabular}

\vspace{0.08cm}
\begin{block}{Why this matters}
These changes removed the most fragile parts of the original design and made the proposal match a realistic implementation path.
\end{block}
\end{frame}

% ============================================================
\begin{frame}{Detectors and Statistical Evaluation}
\footnotesize
\begin{columns}[T]
\begin{column}{0.49\textwidth}
\begin{block}{Submitted version}
\begin{itemize}
\item Detector set centered on SRM+EC plus branch-specific checks.
\item Statistics relied on Wilcoxon, ANOVA, and Cohen's $d$.
\item Trained and training-free detectors were mixed in the main analysis.
\end{itemize}
\end{block}
\end{column}
\begin{column}{0.49\textwidth}
\begin{exampleblock}{Updated version}
\begin{itemize}
\item Primary analysis now uses RS, $\chi^2$, Sample Pairs, and calibration $\chi^2$.
\item ROC-AUC remains the metric, but comparisons now use DeLong for correlated AUCs.
\item Multiple testing is handled explicitly with Bonferroni--Holm.
\item SRNet remains possible, but only as a separate extension.
\end{itemize}
\end{exampleblock}
\end{column}
\end{columns}

\vspace{0.05cm}
\begin{alertblock}{Key methodological correction}
The updated detector/statistics stack fits the actual study structure much better than the original mixed setup.
\end{alertblock}
\end{frame}

% ============================================================
\begin{frame}{Related Work and Reference Corrections}
\footnotesize
\begin{columns}[T]
\begin{column}{0.49\textwidth}
\begin{block}{What was problematic before}
\begin{itemize}
\item The original proposal used a broader but less stable reference set.
\item It treated De et al.\ as the closest prior work.
\item QIM-based framing and some detector references were tied to the dropped methodology.
\end{itemize}
\end{block}
\end{column}
\begin{column}{0.49\textwidth}
\begin{exampleblock}{What the final version does}
\begin{itemize}
\item Related work is now built around the papers that directly support the final design.
\item Gao is used to distinguish coverless/generative steganography from our setting.
\item M\'{e}reur is used as the closest prior comparison on AI-generated carriers.
\item Mallet, Giboulot, and Corvi anchor the cover-source mismatch and synthetic-image forensics argument.
\end{itemize}
\end{exampleblock}
\end{column}
\end{columns}

\vspace{0.05cm}
{\footnotesize\color{umgray}Main change: the literature review became narrower, but more accurate and better aligned with the actual proposal.}
\end{frame}

% ============================================================
\begin{frame}{What Stayed Constant}
\footnotesize
\begin{block}{Core study idea}
\begin{itemize}
\item The project still asks whether carrier origin changes detectability.
\item The main comparison is still real vs.\ ML, with ML-A vs.\ ML-B as a second contrast.
\item The overall scale remains 1,500 images across 3 sources.
\item Payload, encryption, and two embedding branches remain part of the design.
\end{itemize}
\end{block}

\begin{exampleblock}{Practical takeaway for the team}
\begin{itemize}
\item We did \textbf{not} abandon the project direction.
\item We removed risky claims and unrealistic components.
\item The final document is closer to what we can build and defend.
\end{itemize}
\end{exampleblock}
\end{frame}

% ============================================================
\begin{frame}{Net Effect of the Revision}
\footnotesize
\begin{columns}[T]
\begin{column}{0.32\textwidth}
\begin{block}{Improved credibility}
\begin{itemize}
\item fewer hallucination risks
\item cleaner literature positioning
\item more defensible claims
\end{itemize}
\end{block}
\end{column}
\begin{column}{0.32\textwidth}
\begin{block}{Improved methodology}
\begin{itemize}
\item branch-consistent carriers
\item more realistic detectors
\item statistics matched to AUC comparisons
\end{itemize}
\end{block}
\end{column}
\begin{column}{0.32\textwidth}
\begin{block}{Improved communication}
\begin{itemize}
\item less AI-like wording
\item clearer flow
\item tighter explanation of scope and limitations
\end{itemize}
\end{block}
\end{column}
\end{columns}

\vspace{0.08cm}
\begin{alertblock}{Team summary}
The final proposal is best understood as a correction and sharpening of the original submission, not a change of direction.
\end{alertblock}
\end{frame}

\end{document}
